\subsection*{Q and A}

\walkfig{a1-intro}{The Q and A tab opens on a single promise about vocabulary.}

\Stu{This tab is the skeptic's index. Fourteen questions, all closed at load. The rule in that intro line is the rule the whole app runs on: every answer either jumps into a tab where the claim can be recomputed on the spot, or it carries a name and a year --- and the three words \emph{theorem}, \emph{computed live}, and \emph{analogy} never trade places.}

\Kle{Fourteen questions a skeptical manifold theorist asks. Written by you, in my voice, before I got here.}

\Stu{Yes. That is the obvious objection and I will not dress it up: I chose the questions, so I chose the ones I could answer. What I can defend is that none of the fourteen has a soft answer inside it. Three of them are failure lists.}

\Kle{Fine. I will supply my own questions as we go and we will see which ones you anticipated. Open the first.}

\walkfig{a2-q1}{Question one: why should a prime resemble a knot at all.}

\Stu{Two theorems and one leap, kept strictly apart. Theorem: $\mathrm{Gal}(\bar{\mathbb{F}}_p/\mathbb{F}_p) \cong \hat{\mathbb{Z}}$, topologically generated by Frobenius --- the same profinite group as the profinite $\pi_1$ of the circle, so $\mathrm{Spec}\,\mathbb{F}_p$ is a $K(\hat{\mathbb{Z}}, 1)$ exactly as $S^1$ is a $K(\mathbb{Z},1)$. Theorem: $\mathrm{Spec}\,\mathbb{Z}$ behaves three-dimensionally --- etale cohomological dimension 3, with Artin--Verdier duality playing its Poincare duality, Artin--Verdier 1964 and Mazur 1973. And then the leap, in bold and labelled Analogy: a circle sitting in a 3-space is a knot. Mazur's observation, unpublished notes around 1963 to 1964.}

\Kle{Read me the sentence after the leap.}

\Stu{``Asserts no embedding and no isotopy; that leap \emph{is} the dictionary.''}

\Kle{Good. That is the sentence that matters and I am glad it is on the page rather than in your mouth. Because the honest statement of the situation is that you have a $K(\pi,1)$ on one side, a duality of the right formal shape on the other, and nothing whatsoever connecting them --- no map, no embedding of $\mathrm{Spec}\,\mathbb{F}_p$ into $\mathrm{Spec}\,\mathbb{Z}$ that behaves like an embedding of $S^1$ into $S^3$. There is a map of schemes, and it is not an embedding in any sense I would use the word.}

\Stu{Agreed, and the app never claims one. What survives the leap is that both duality statements are top-degree 3, so the pairings have matching arities --- which is what lets a linking number be defined on both sides at all.}

\Kle{That is the right defence. Say it that way. ``Matching arities'' is a claim I can check; ``a prime is a knot'' is a slogan.}

\begin{knote}
Q\&A tab. Q1 is honest where it counts: theorem / theorem / analogy, in bold, in that order. He conceded instantly that there is no embedding. The real content is that both dualities are top-degree 3, so the pairings have the same arity --- which is what a linking number needs. Tell him to promote that sentence into the answer; it is currently only in his voice, not on the page.
\end{knote}

\walkfig{a3-q9}{Question nine open: are the two zeta functions the same object?}

\Stu{The answer starts with the word ``No,'' which is the point. The Lefschetz zeta of a monodromy and the Dedekind zeta of a number field are related by analogy, not equality --- and the Zeta tab says so in a note above the two columns, before you can object. What the columns actually share is a mechanism: one sum resummed as a product two ways. Iterates, eigenvalues, orbits on the left; ideals, characters, primes on the right, and every instance reduces to $-\log(1-x) = \sum x^n/n$.}

\Kle{Where does equality ever get asserted, then?}

\Stu{Only within a column, never across. Two up-to-units identities are theorems, each inside one side of the mirror: $\tau_L \doteq \Delta_L$, Milnor 1962 and Turaev 1986, and the Main Conjecture $f \doteq L_p$, Mazur--Wiles 1984, with $\ell = 2$ by Greither 1992. And the one numerical equality on the tab is $\zeta_{\mathbb{Q}(\sqrt{5})}(2)$ computed three ways --- ideal sum, Euler product, character product --- converging to the same decimals, live.}

\Kle{Three ways to compute one number is not three theorems, it is one theorem and two implementations. But I take the point of the display, which is that you are willing to be caught.}

\walkbeat{He reaches over and opens the next question down himself.}

\Kle{This is the one I actually wanted. ``What exactly is computed live, and what am I being asked to take on faith.'' Read me the faith list.}

\Stu{Cited, not computed: fiberedness, Stallings 1978; $\bar\mu(123) = \mathrm{lk}(\gamma, K_3)$, Cochran 1990 and Mellor--Melvin 2003; the model longitudes, Milnor 1957; Redei's ramification statement, 1939; Amano's independence theorem, Tohoku 2014; Milnor--Turaev, Mazur--Wiles, Greither. Computed: every Legendre symbol by Euler's criterion, Tonelli--Shanks square roots, Redei symbols with the normalization enforced, Magnus coefficients, matrix orders by literal powering, Bernoulli numbers mod $p$.}

\Kle{And the flagged gap.}

\Stu{One, on the Quadruples tab: $\bar\mu(1234)$ is computed for Milnor's model $B_4$, not for the drawn diagram. The Wirtinger presentation is not implemented and the honesty note under the picture says so.}

\Kle{Here is my complaint about this entire tab, and I want it recorded. Not one number appears on it. This is the tab whose whole subject is what is computed live, and it computes nothing --- it is fourteen paragraphs of prose with buttons. Every claim's support is one click away and none of it is visible. If I read only this tab I would have to take the word ``computed'' entirely on trust, which is precisely the posture the tab exists to refuse.}

\Stu{That is fair and I have no answer to it. The tab was built as an index, and an index that indexes computation is not computation.}

\Kle{Put one live number on it. One. Recompute $(13/61)$ in the margin of question one and the tab stops being a brochure.}

\walkfig{a4-breakage}{The last question: where the dictionary breaks, five fractures.}

\Stu{Five, each flagged in the app at its point of use. One: the archimedean place --- arithmetic carries $S \cup \{\infty\}$ bookkeeping everywhere, and the whyFS box on the Dictionary tab shows it biting at $\ell = 2$, where the totally real tower drops to $\mathbb{Z}/4 \times \mathbb{Z}/2$; the knot side has no mirror for $\infty$ past the $p \equiv 1 \pmod 4$ orientability convention. Two: no arithmetic hyperbolic geometry --- the figure-eight complement carries a volume and dilatation $\varphi^2$, and no prime carries anything like a volume. Three: the $B_4$ gap. Four: the ring mismatch --- $\Delta_L$ lives over $\mathbb{Z}[t_1^{\pm 1}, \ldots, t_r^{\pm 1}]$ with deck group $\mathbb{Z}^r$, and over $\mathbb{Q}$ no $\mathbb{Z}^r$-cover exists. Five: it is not a functor.}

\Kle{Stop at four. Say the fourth one again, slowly, because it is the only one on the list that is a theorem rather than a shrug.}

\Stu{Over $\mathbb{Q}$ there is no $\mathbb{Z}^r$ deck group unramified outside the $p_i$. The whyFS box computes the $f_S$ deck group finite --- $\mathbb{Z}/4 \times \mathbb{Z}/4$ for $S = \{13, 61\}$ at $\ell = 2$. And the one $\mathbb{Z}_p$-tower $\mathbb{Q}$ does have, the cyclotomic one, ramifies at $p$ itself.}

\Kle{So the two polynomials do not merely happen to live over different rings --- they are \emph{forced} apart, by finiteness on one side. That is a much stronger and much better statement than ``the fit is imperfect,'' and it is buried as item four of five in a paragraph. Fractures one, two, three and five are honest caveats. Four is a theorem about why the analogy cannot be upgraded. Lead with it.}

\Stu{Noted. It is currently ranked by where it appears in the app, not by weight.}

\Kle{And number five --- ``it is not a functor.'' Who is that for? A manifold theorist hears that and asks which category you failed to construct. A number theorist may not hear anything at all. It is the correct statement, but it is the only one of the five that is a claim about mathematics rather than about your app, and it gets one sentence.}

\begin{knote}
Fracture (4) is the real one and it is buried: over $\mathbb{Q}$ no $\mathbb{Z}^r$-cover unramified outside $S$ exists, $f_S$ deck group finite ($\mathbb{Z}/4 \times \mathbb{Z}/4$ for $S = \{13,61\}$, $\ell = 2$), and the genuine $\mathbb{Z}_p$-tower ramifies at $p$. So $\Delta_L$ and the Iwasawa polynomial are forced over different rings --- not a mismatch of taste, a finiteness theorem. This is the sharpest thing on the tab. It is item four of five in a run-on paragraph. Also: this tab computes nothing. Tab about live computation, zero live computation.
\end{knote}

\walkfig{a5-list}{All fourteen questions, closed --- the tab as it loads.}

\Stu{This is the resting state. Everything shut, so the tab reads as a contents page rather than a wall.}

\Kle{Then let me review it as a contents page. They are not numbered. You told me in the first line there are fourteen, so I have to count them myself to find the fourteenth, and there is no way to open them all at once --- if I want to read this tab straight through I click fourteen times. That is a small thing and it is the difference between a document and a toy.}

\Stu{There is no expand-all control. Correct.}

\Kle{Second thing, and this one is not small. Count the buttons that jump out of this tab.}

\Stu{Twenty, across six destinations --- Dictionary, Pairs, Triples, Ladder, Quadruples, Zeta.}

\Kle{Six. There are seven other tabs. Not one of these fourteen answers sends me to the Experimental tab, which is where you keep the things you are least sure about. Your skeptic's index routes the skeptic away from the frontier. If question thirteen --- the five-component one --- had a button to it, I would have gone there under my own steam instead of because you walked me.}

\Stu{That is a genuine omission. Question thirteen ends by saying Amano's quadruple is the frontier, and the tab that argues about frontiers is the one it does not link.}

\Kle{One more. Question six --- ``everything in this app is mod 2.'' Its answer uses the phrase ``the $f_S$ deck group'' twice and never says what $f_S$ is. I can guess: the maximal pro-$\ell$ extension unramified outside $S$. But I am guessing, and I am the audience.}

\begin{knote}
Contents page: unnumbered, no expand-all, fourteen clicks to read. Twenty jump buttons, six destinations, none to Experimental --- the index avoids the frontier. $f_S$ used twice in Q6, defined nowhere on the tab. Overall: the honesty vocabulary holds up across all fourteen. I could not find a place where \emph{theorem} was doing work that only \emph{analogy} had earned. That is rarer than it sounds.
\end{knote}

\subsection*{Experimental}

\walkfig{x1-banner}{The frontier banner: nothing here is load-bearing.}

\Stu{The banner sets the rules for the whole tab. Nothing here is load-bearing for the other seven. Every claim below is marked one of four ways --- \emph{theorem}, \emph{computed live}, \emph{empirical} meaning machine-checked but not derived, or \emph{proposal / conjectural}. Two threads: what the zeta function counts when you change the generator, and what happens to the dictionary at $\ell = 3$.}

\Kle{The four-way labelling is the right idea and I want to be clear that I approve of it. Most people have three tiers at best. ``Empirical, machine-checked but not derived'' is a category most mathematicians refuse to admit exists in writing, and you have given it a name and a badge.}

\Stu{It exists whether or not it is admitted.}

\Kle{It does. Now --- look at how it is set. That paragraph is the constitution of the tab and it is rendered in small grey centred italics. It is the least legible block on the page and it governs everything below it. Whoever chose that style was thinking ``modest,'' and modesty in typography reads as ``skippable.''}

\Stu{The style class is shared with every honesty box in the app --- point-eight-five rem, muted, italic, centred. It is applied to three blocks on this tab and two of them are the longest paragraphs here.}

\Kle{Then it is a systematic error, not a slip. Your most important prose is your least readable prose.}

\walkfig{x2-generators}{The two ramified generators, side by side: the trefoil worked live against the Iwasawa twin.}

\Stu{Left column, the knot story. Feed the Lefschetz formula the monodromy $\varphi_*$ of a fibered knot acting on $H_1$ of its fiber --- a local, ramified generator. Milnor 1968: you get the Alexander polynomial. The trefoil is fibered over $S^1$ with fiber a once-punctured torus and monodromy the $SL_2(\mathbb{Z})$ matrix $[[1,1],[-1,0]]$, characteristic polynomial $\Delta(t) = t^2 - t + 1$. That polynomial is cited. Everything in the table is computed from it here: $\#H_1(M_n) = \prod_{\zeta^n = 1}|\Delta(\zeta)|$, giving 3, 4, 3, 1 at $n = 2, 3, 4, 5$, and infinite at $n = 6$.}

\Kle{Three, four, three, one, infinite.}

\walkbeat{He stops the walkthrough and looks at the column of integers for several seconds.}

\Kle{Those are the Brieskorn spheres. $\Sigma(2,3,n)$. Your $n = 2$ row is the lens space $L(3,1)$, first homology $\mathbb{Z}/3$. Your $n = 3$ is $\Sigma(2,3,3)$ with $\mathbb{Z}/2 \oplus \mathbb{Z}/2$, order four. And your $n = 5$ row, the one you have printed as the integer 1 --- that is the Poincare homology sphere. $\Sigma(2,3,5)$. The most famous 3-manifold in my subject is sitting in your fourth table row labelled ``1'' and you have not named it.}

\Stu{I had not made that identification. The app computes the order and stops.}

\Kle{And $n = 6$ infinite, because $\Delta = \Phi_6$ vanishes at a primitive sixth root --- $\Sigma(2,3,6)$, which is not even hyperbolic, it is the Nil one. Your table is a complete tour of the trefoil's branched covers and it presents itself as a column of arithmetic.}

\Stu{Which brings up something you should say back to me, because I think you are about to.}

\Kle{I am. Nowhere on this tab do you write the word ``branched.'' You say ``the whole tower of cyclic covers,'' and $M_n$ is never defined. But that product formula is the formula for the $n$-fold cyclic \emph{branched} cover of $S^3$ branched over $K$. The unbranched cyclic covers have different homology --- the infinite cyclic cover of the trefoil has $H_1 = \mathbb{Z}^2$ as a group, not order 3. A reader who takes your word ``cyclic covers'' literally will compute the wrong thing.}

\Stu{You are right, and I have checked: ``branched'' appears exactly once in the entire file, in an unrelated Dictionary tooltip about ramification. On this tab it is missing and $M_n$ is undefined. That is a straightforward defect.}

\Kle{It is the good kind of defect --- the numbers are right and the word is missing. Now the right column.}

\Stu{The same formula with the other ramified generator: $\gamma$, the topological generator of a $\mathbb{Z}_p$-extension acting on the Iwasawa module $H_\infty = \varprojlim \mathrm{Cl}(k_n) \otimes \mathbb{Z}_p$. Iwasawa: you get $I_p(T) = \det(T \cdot \mathrm{id} - (\gamma - 1))$ up to $p^\mu$, with $\#H_n = p^{\lambda n + \mu p^n + \nu}$. And the quantity that matches on the topological side is not the order but $v_p(\#H_1(M_{p^n}))$ --- the third column at left. For the trefoil at $p = 2$ that column is identically zero, because 3 and 3 are odd, so $\lambda = \mu = \nu = 0$. The asymptotic in its degenerate case.}

\Kle{You have matched two things by exhibiting an instance where both are zero. That is not a match, that is a coincidence of vanishing.}

\Stu{Agreed, and I would not call it evidence. What it does is force the right statement: the correspondence has to be with $v_p$, not with the integer, because $\#H_1$ need not be a $p$-power at all --- your $n = 3$ row is 4, and 4 is not a power of 2 in any way that would help if $p = 3$. The degenerate case is what makes the shape of the claim visible. It is not what supports it.}

\Kle{Accept that. Now, this column has one genuinely important sentence in it and it is not visible. Where is the statement that $\gamma$ is \emph{not} analogous to Frobenius?}

\Stu{In a hover tooltip. The words ``not to Frobenius'' are underlined and the explanation appears on hover: $\gamma$ generates the ramified tower itself, Frobenius generates an unramified local extension --- the loop around a puncture disjoint from the branch locus. Conflating them is the commonest error in reading the dictionary.}

\Kle{Your own tooltip calls it the commonest error, and you have put the correction behind a mouse. Also --- look at the shape of this screen. Your left column runs to the bottom of the viewport and your right column stops halfway, so there is a column of empty page next to your most delicate paragraph. It reads as though the right side ran out of things to say.}

\begin{knote}
Best beat of the tab and he did not see it: his trefoil table is $\Sigma(2,3,n)$ for $n = 2,3,4,5,6$ --- $L(3,1)$, $\mathbb{Z}/2^2$, $\mathbb{Z}/3$, the \emph{Poincare sphere}, and Nil. He prints ``1'' where $\Sigma(2,3,5)$ belongs. Name them; the whole point of this app is that the two columns illuminate each other, and here the topological column has famous names he is not spending.

Defect: ``branched'' never appears, $M_n$ never defined; the product formula is for the branched covers. Numbers right, word missing.

$\lambda = \mu = \nu = 0$ is agreement by mutual vanishing. He conceded it unprompted and gave the better reason for it (matching must be with $v_p$, not the order, since $\#H_1$ need not be a $p$-power). Correct.

The load-bearing sentence of the tab --- $\gamma$ is not Frobenius --- is in a hover tooltip.
\end{knote}

\walkfig{x3-pivot}{The pivot box: change the generator, change the subject.}

\Stu{This is the hinge of the tab. Feed the Lefschetz formula the ramified generator, $\varphi_*$ or $\gamma$, and it returns $\Delta_K$ and $I_p$ --- knot-type, depth two, blind to $\bar\mu(123)$. Feed it the unramified generator, Frobenius --- the loop around a puncture disjoint from the branch locus --- and the same formula is an Artin $L$-function, which is where the Redei symbol and every higher-linking invariant in this app already live.}

\Kle{Say the second half of the box.}

\Stu{The app read back against itself: Pairs and Triples pair against Frobenius, so they are the link column, unramified. The Zeta tab's figure-eight monodromy card and the knotted-prime card's Iwasawa polynomial use $\varphi_*$ and $\gamma$, so they are the knot column, ramified. Same formula, two seats --- a distinction the other tabs use without naming.}

\Kle{That is the best paragraph in the app and it is on the tab you told me was not load-bearing.}

\Stu{I know.}

\Kle{No, I want you to hear the size of it. You spent seven tabs computing linking numbers and Iwasawa invariants side by side, and the actual reason they are different kinds of object --- one generator ramified, one not --- is quarantined on the frontier tab under a banner saying ``nothing here is load-bearing.'' It is load-bearing. It is the load. Every reader who leaves after the Zeta tab leaves thinking $\Delta_K$ and $\bar\mu(123)$ are the same flavour of invariant, and they are not: one is depth two and one is not.}

\Stu{The banner is about epistemic status, not about importance --- nothing here is \emph{relied on} by the other tabs.}

\Kle{Then your banner is measuring the wrong axis, and it is costing you your best idea. Move this box to the Dictionary tab and label it theorem-plus-definition, because that is what it is. The conjectural material can stay out here.}

\walkfig{x3b-cracks}{The two honest cracks, stated because they are load-bearing.}

\Stu{Two, and both are conceded rather than defended. First: the indeterminacy does not match. Milnor's torsion $\doteq \Delta_K$ holds up to $\mathbb{Z}[t,t^{-1}]^\times = \{\pm t^n\}$, a tiny group. The Main Conjecture's ``up to units'' is up to $\Lambda^\times$ where $\Lambda = \mathbb{Z}_p[[T]]$ --- every power series with unit constant term. Enormously larger. The triangle repeats; its slack does not; and the note says outright that it does not know a reason it should.}

\Kle{That is the correct thing to be worried about and almost nobody writes it down. $\{\pm t^n\}$ is countable and discrete. $\Lambda^\times$ is a pro-$p$ group of infinite rank. An identity ``up to $\Lambda^\times$'' pins a principal ideal, and an identity up to $\{\pm t^n\}$ pins a polynomial to sign and shift. Those are not the same strength of statement by any measure, and if the analogy were structural rather than formal, the slack ought to correspond too.}

\Stu{That is exactly the sentence the box is trying to say and yours is better.}

\Kle{Second crack.}

\Stu{The depth-three statement is a proposal, not a theorem. The reading offered is that the order of vanishing of the Artin $L$-function of the tower generated by the arithmetic Seifert surface $\sqrt{\alpha}$ is the number-theoretic analogue of Cochran's derived invariant, and that the Redei symbol sits inside such an order of vanishing.}

\Kle{``One rung below the machine-checked depth-four case.'' What is the depth-four case?}

\Stu{The $d$-bit law in the box below.}

\Kle{Then say so. The phrase ``depth-four'' appears twice on this tab and is defined neither time, and the thing it refers to is called something else entirely where it appears --- I have just read the next box and it is called ``the one machine-checked instance'' and ``a governing bit.'' You have three names for one object on one screen. And ``depth two'' you did define, in the left column, so a reader has every reason to think depth-four means four link components. Does it?}

\Stu{No. It does not mean four components.}

\Kle{Then that is a term doing real work with no definition and a misleading neighbour. Fix it or drop it.}

\walkfig{x4-redei}{A second Borromean prime triple, computed live: $(13, 17, 53)$.}

\Stu{This is a live computation, not a citation. Pairwise: $(13/17) = +1$, $(13/53) = +1$, $(17/53) = +1$ --- all three, the arithmetic unlink, and all three primes $\equiv 1 \pmod 4$. Redei witness $(x,y,z) = (-15, 4, 1)$: the identity $x^2 - 13y^2 - 17z^2$ evaluates to 0 on screen, $y$ is even, $x - y \equiv 1 \pmod 4$. Conjugates of $\alpha$ mod 53 come out 45 and 31, both Legendre $-1$, so $[13,17,53] = -1$.}

\Kle{Pairwise unlinked, triply linked, on a triple I have not seen before. Good --- that is the answer to the cherry-picking question your Q and A tab gave in prose, given here in numbers instead. Why is $x$ negative?}

\Stu{Some pairs have no normalized certificate with $x$ positive. The engine accepts any sign on the certificate coordinates; only the primes must be positive.}

\Kle{And it is displayed as ``-15'' with a hyphen while your symbol value four lines later is ``$-1$'' with a proper minus sign. On the same card. I mention it because you have clearly been careful about typography everywhere else, so this one is going to look like the number came from somewhere else --- which, in a document about what is computed live versus what is quoted, is exactly the wrong impression to give.}

\Stu{It is a raw join of the certificate array. It should go through the same formatter as everything else.}

\Kle{Now the sentence underneath: it is $\mathrm{Frob}_{53}$ paired against $\sqrt{\alpha}$ that decides it. Unramified generator, doing link theory. That is the pivot box being cashed out one screen later, and it is the first time on this tab that the abstract distinction has produced a number. Say that when you present this --- the reader should see the pivot pay.}

\walkfig{x5-dbit}{The $d$-bit box: the empirical law, this app's own pair, and the normalization caution.}

\Stu{Three blocks, three different epistemic statuses, and they are labelled. Block one is empirical and is not computed here: the laboratory reports, on 8091 pairs, that a governing bit $d\{p,q\} \in \{1,3\}$ defined by the sign of the Redei witness $\alpha$ --- $d = 1$ when $F_{p,q} = \mathbb{Q}(\sqrt{p}, \sqrt{q}, \sqrt{\alpha})$ is totally real, $d = 3$ when totally imaginary --- matches an a priori independent analytic invariant, the order of vanishing of the Artin $L$-function: $d = 1$ if and only if $\mathrm{ord}_{s=0} L(s, \rho_{p,q}) = 2$, and $d = 3$ if and only if the order is 0. The equality of two independently defined quantities is the claim. No derivation is asserted.}

\Kle{Why one and three? Not zero and one, not plus and minus. What are you indexing?}

\Stu{The box does not say, and I should be honest that I cannot reconstruct it from the app. The labels come from the source note.}

\Kle{Then on this page they are decoration. A bit with two values called 1 and 3 invites me to think there is a $d = 2$ somewhere, and there is not. Same for ``a Stark rank'' in parentheses --- that is three words standing in for a theory, and this tab has already shown it is willing to spend a paragraph on things that matter.}

\Stu{Block two is computed live: our own pair. $(13, 61)$ with the shipped normalized certificate $(23, 6, 1)$ gives $\alpha = 23 + 6\sqrt{13} = 44.633$ and $\bar\alpha = 1.367$. Both positive, so totally positive, so $k_{13,61}$ is totally real and $d\{13,61\} = 1$.}

\Kle{And block three is the one you should lead with.}

\Stu{Block three is a caution the app found while checking, and states rather than hides. The sign is only well defined once Redei's normalization is imposed. For $(13,17)$, every one of the 10 normalized witnesses with $z \le 12$ gives $\alpha$ totally negative --- 10 negative, 0 positive, computed at load --- so $d = 3$. But the illustrative witness $(9,1,2)$ quoted in the source note is primitive and non-degenerate and \emph{not} Redei-normalized, because $y$ is odd; our engine throws on it. And it gives $\alpha = 12.606$, $\bar\alpha = 5.394$, totally positive, so $d = 1$. Same pair, opposite bit.}

\Kle{You found a counterexample to your own source's worked example and you shipped it in the app rather than quietly fixing the example.}

\Stu{The invariant is fine. The definition needs the normalization clause and mere primitivity is not enough. That was worth more than the example was.}

\Kle{It is the single most creditable thing on this tab and I will say so plainly. Now let me break it. ``Every one of the 10 normalized witnesses with $z \le 12$.'' What is the bound on $y$?}

\Stu{The scan runs $y$ from $-120$ to $120$.}

\Kle{Which is not stated. You have declared one of two bounds and phrased the result as ``every one.'' That is the shape of an exhaustive claim over a range you have only half-described. It does not matter that the conclusion is almost certainly right --- on a tab whose entire subject is labelling what is checked and what is not, an undeclared search bound is exactly the wrong omission.}

\Stu{Conceded without qualification. The prose should carry both bounds or neither.}

\begin{knote}
$d$-bit box: three blocks, three statuses, labelled --- empirical (8091 pairs, quoted not computed), computed live (13, 61 totally positive, $d = 1$), and the caution. The caution is the best work in the app: they checked their own source's illustrative witness $(9,1,2)$, found it unnormalized ($y$ odd), found it flips the bit --- same pair, opposite bit --- and published it. That is a laboratory, not a demo.

Against it: search bound half-declared ($z \le 12$ stated, $|y| \le 120$ not). ``Every one of the 10'' is an exhaustiveness claim over an undescribed box. And $d \in \{1,3\}$ is never explained; ``a Stark rank'' is three words for a theory.
\end{knote}

\walkfig{x6-cubictable}{Ten rows: what changes at $\ell = 3$.}

\Stu{Second thread. Everything else in the app is mod 2 --- values in $\{\pm 1\}$, double covers, deck groups $D_4$ and $N_4(\mathbb{F}_2)$. The framework is not restricted to $\ell = 2$. This table is the audit: ten rows, what the $\ell = 2$ column says and what the $\ell = 3$ column says. Value group goes from $\{\pm 1\} = \mu_2$ to $\mu_3 = \{1, \omega, \omega^2\}$; base field from $\mathbb{Q}$ to $k = \mathbb{Q}(\zeta_3) = \mathbb{Q}(\sqrt{-3})$; orientability convention from $p \equiv 1 \pmod 4$ to $p \equiv 1 \pmod 3$, which is what makes $p$ split in $k$; normalization from Redei's $x - y \equiv 1 \pmod 4$ to \emph{primary}, $\pi \equiv 2 \pmod 3$ in $\mathbb{Z}[\omega]$; double cover to 3-fold cyclic cover; quadratic reciprocity to cubic; $N_3(\mathbb{F}_2) \cong D_4$ of order 8 to $N_3(\mathbb{F}_3)$ of order 27 with centre $\mathbb{Z}/3$.}

\Kle{Row eight. ``Chebotarev share per class at the centre --- one half each, one third each.'' Your Q and A tab told me the Ladder's Chebotarev densities were one eighth. Which is it?}

\Stu{Both, and the row is compressed to the point of being misleading. On the Ladder the density of the identity in $D_4$ is $1/8$ and the density of the central corner matrix is $1/8$ --- densities in the whole group. This row is conditional: given that Frobenius has landed in the centre, which for $D_4$ is $\mathbb{Z}/2$ of order two, each of the two elements takes half. For $N_3(\mathbb{F}_3)$ the centre is $\mathbb{Z}/3$ and each takes a third. Same theorem, different denominator.}

\Kle{And the word ``conditional'' is not in the row.}

\Stu{No.}

\Kle{Then a reader who has just come from the Ladder tab sees $1/8$ turn into $1/2$ and has to reconstruct your denominator. Two extra words fixes it. Second point --- the last row calls it ``the $d$-defect,'' and the box we just read calls it ``a governing bit $d\{p,q\}$,'' and the crack above that calls the same territory ``depth-four.'' Three vocabularies for one thread on one tab.}

\Stu{Fair. The table row was written to be scannable and it drifted.}

\Kle{The last two rows are in red and say ``none here'' and ``no real place.'' That is the right instinct --- the failures are in the table rather than in a footnote, so I cannot read the comparison without reading the failures. Keep that.}

\walkfig{x7-mod3topo}{The topology mod 3: three dial positions, $|{\rm deck}| = 3 = $ loops times winding.}

\Stu{Nothing changes in kind, only the modulus. In the 3-fold cyclic cover of $S^3 \setminus K_1$, a loop with monodromy of order $d$ has $3/d$ preimage components each winding $d$ times, so the dial is $\mathrm{lk} \bmod 3$ with three positions instead of two. $\mathrm{lk} \equiv 0$: order 1, three loops winding 1, orbit factor $(1-t)^{-3}$. Otherwise: order 3, one loop winding 3, factor $(1-t^3)^{-1}$. And $3 = \text{loops} \times \text{winding}$ in every row, computed.}

\Kle{This is the row of the dictionary I have no complaint about, because it is just covering space theory and covering space theory does not care what $\ell$ is. Index times sheets. It would be strange if this broke.}

\Stu{Which is why it is the part I am most confident of and least impressed by.}

\Kle{Correct posture. Small thing: you have five rows for three positions. Rows for $\mathrm{lk} = 3$ and $\mathrm{lk} = 4$ are literally rows 0 and 1 again. I assume that is deliberate, to show periodicity --- but nothing on the table says so, so it looks at first glance like the dial has five positions, three sentences after you told me it has three.}

\Stu{It is a periodicity demonstration with no label saying it is one.}

\Kle{And the note at the bottom --- $\bar\mu(123) \bmod 3$ landing at the corner of $N_3(\mathbb{F}_3)$, ``described not computed,'' no mod-3 Magnus pipeline, so no number asserted. Good. That is the sentence I would have demanded if it were not there.}

\walkfig{x8-cubicworked}{The Eisenstein arithmetic, computed live: primary associates and the cubic residue symbol both ways.}

\Stu{Live, at load. Both 7 and 61 are $\equiv 1 \pmod 3$, so both split in $\mathbb{Z}[\omega]$. The app finds a prime above each --- $-3 - 2\omega$ above 7, $-9 - 5\omega$ above 61 --- and then normalizes each to its primary associate by multiplying through the six units until $\pi \equiv 2$ and the $\omega$-coefficient $\equiv 0 \bmod 3$. That gives $\pi_7 = -1 - 3\omega$ and $\pi_{61} = -4 - 9\omega$, and the norms are recomputed: 7 and 61. Then the symbol $\chi_\pi(\alpha) \equiv \alpha^{(N(\pi)-1)/3} \pmod \pi$, valued in $\mu_3$, computed in the residue field, with $\omega \mapsto 2 \bmod 7$ --- and 2 does satisfy $x^2 + x + 1 \equiv 0 \bmod 7$, which the app checks. Result: $\chi_{\pi_7}(\pi_{61}) = 1$ and $\chi_{\pi_{61}}(\pi_7) = 1$. Equal. Cubic reciprocity.}

\Kle{Both of them are 1.}

\Stu{Yes.}

\Kle{So your worked instance of cubic reciprocity is the identity equals the identity. You have verified $1 = 1$ and put a green tick on it. If your implementation of $\chi$ returned the constant 1 for every input, this card would look exactly the same.}

\Stu{That is right and it is a real weakness of the worked example. The evidence that the implementation does anything lives one paragraph down, in the census, where the values are not all trivial.}

\Kle{Then swap them, or add a second row to this card with a nontrivial pair. Choosing a degenerate worked example is a strange thing to do on the one tab where you have made a virtue of finding your own degenerate cases.}

\Stu{The one thing I would keep is the gloss under it: both trivial means 7 and 61 are cubically unlinked --- 61 is a cube mod 7 and 7 is a cube mod 61 --- the mod-3 analogue of Legendre symbol $+1$. That sentence is what makes the number mean something rather than just be equal.}

\Kle{Keep the sentence, change the primes. And the closing line is right: symmetry is what earns the name linking number. That is a real principle and it is the reason the arithmetic side is allowed to say ``linking'' at all --- $\mathrm{lk}(K_1, K_2) = \mathrm{lk}(K_2, K_1)$ is not a convention, it is a theorem, and so is this.}

\walkfig{x9-census}{The 171-pair cubic reciprocity census, computed at load.}

\Stu{All 19 primes $\equiv 1 \pmod 3$ below 200, primary-normalized, and all 171 unordered pairs. Reciprocity $\chi_\pi(\theta) = \chi_\theta(\pi)$ holds 171 out of 171. Value distribution $1 : \omega : \omega^2 = 61 : 57 : 53$, which is 35.7 / 33.3 / 31.0 percent, against the Chebotarev prediction of a third each.}

\Kle{The 171 out of 171 is not evidence of anything mathematical. Eisenstein proved cubic reciprocity. What you have tested is your implementation.}

\Stu{Correct, and the card does say Eisenstein is cited. But it does not say that the 171/171 is a unit test rather than a result, and it should.}

\Kle{Now the distribution, which is the part that is actually a measurement. 61, 57, 53 out of 171. Expected 57 each. What is your standard deviation?}

\Stu{For $n = 171$ at probability one third, the variance is $171 \times (1/3)(2/3) = 38$, so about 6.2.}

\Kle{So your deviations are plus four, zero, minus four --- about two thirds of one standard deviation. Which is to say: perfectly consistent with a third each, and also perfectly consistent with a good many other things, because 171 pairs is nothing. And you have printed 35.7 and 31.0 to one decimal place with no error bar, which invites a reader to see a trend in noise. Either print the standard deviation next to it, or do not print the third significant figure.}

\Stu{That is right. The percentages are more precise than the sample can support.}

\Kle{The closing sentence is honest, at least: your laboratory's real census runs 120 thousand triple symbols and this app cannot compute a single one, because there is no mod-3 triple-symbol construction here. Say that number louder --- the gap between 171 pairs and 120 thousand triples is the whole distance between this tab and the research programme behind it.}

\begin{knote}
Cubic side: primary normalization and $\chi_\pi$ are genuinely computed --- $\pi_7 = -1 - 3\omega$, $\pi_{61} = -4 - 9\omega$, norms recomputed, $\omega \mapsto 2 \bmod 7$ verified against $x^2+x+1$. Solid.

But the worked reciprocity instance is $1 = 1$. Degenerate. A constant function would pass it. The nontrivial evidence is in the census one paragraph below and he knows it.

Census: 171/171 reciprocity tests the code, not the mathematics (Eisenstein proved it) --- not said. Distribution 61:57:53, expected 57, sd $\approx 6.2$; deviations $\approx 0.65$ sd. Consistent, and meaningless at $n = 171$. Printed to one decimal with no error bar.

Covering-space row (mod-3 topology) is unimpeachable and he said so himself --- index times sheets does not care about $\ell$.
\end{knote}

\walkfig{x10-gaps}{The $\ell = 3$ gap list: no Redei construction, no real place, the honesty ladder.}

\Stu{Three gaps. One: no Redei construction. Redei's $\alpha$ is quadratic by design --- it comes from the conic $x^2 - p_1y^2 - p_2z^2 = 0$. The $\ell = 3$ analogue needs a degree-3 Kummer relative-norm solver, which this app does not have, so the mod-3 triple symbol is described here and never computed. Two: the defect has nowhere at infinity to live. At $\ell = 2$ the governing bit is a sign --- a real place, totally real against totally imaginary. At $\ell = 3$ the base field is $\mathbb{Q}(\sqrt{-3})$, imaginary quadratic, zero real embeddings. There is no real place for a sign to sit in, so the conjecture is that the defect migrates to the ramified prime $\lambda = (\sqrt{-3})$ above 3 --- a $\lambda$-adic $d_\lambda$ replacing an archimedean sign. Labelled conjectural, and it is the crux of the programme.}

\Kle{Stop. Gap two says ``zero real embeddings'' and then, in parentheses, ``computed at right.'' Show me where at right that is computed.}

\walkbeat{The student scrolls up to the right-hand column and stops.}

\Stu{It is not. The right-hand column computes primary associates, cubic residue symbols and reciprocity. It never computes a signature. $\mathbb{Q}(\sqrt{-3})$ having no real embeddings is stated in the table row above, not computed anywhere.}

\Kle{So on the tab whose entire discipline is the four-way label --- theorem, computed live, empirical, conjectural --- you have a parenthesis reading ``computed at right'' pointing at a computation that does not exist. That is the one error I would call serious, because it is not a gap in the mathematics, it is a false label on the honesty system itself. Everything else on this tab I can check by reading. That one I could only catch by scrolling.}

\Stu{I have no defence and I am not going to construct one. It is a wrong word in the one place the app cannot afford a wrong word.}

\Kle{Which is, in its way, a compliment to the rest of it --- I had to look at all eleven boxes to find one. Read me the third gap.}

\Stu{The honesty ladder. $\ell = 2$ is proven: Legendre, Redei, Amano, and the one machine-checked depth-four law. The $\ell = 3$ migration to $d_\lambda$ is conjectural. The $p$-adic founding conjecture behind both is speculative. This tab exists to keep those three tiers visibly apart.}

\Kle{Three tiers, and the sentence naming them is set in small grey centred italics at the bottom of a wall of small grey centred italics. I said this at the top of the tab and I will say it once more because it is the same defect three times: your gap lists are your best writing and your worst typography. A reader skims grey italic. You have put the load in the place the eye is trained to skip.}

\begin{knote}
Serious: gap (ii) says $\mathbb{Q}(\sqrt{-3})$ has zero real embeddings ``(computed at right)'' --- nothing at right computes a signature. A false ``computed'' label on the tab whose whole subject is what is computed. Not a mathematical error; a labelling error, which here is worse.

Otherwise the gap list is the right list: no degree-3 Kummer solver so no mod-3 triple symbol, ever, in this app --- said outright. The $d_\lambda$ migration conjecture is clearly marked. The three-tier ladder (proven / conjectural / speculative) is exactly the structure this material needs.

Typography, third time: the three most important paragraphs on the tab are the three least readable. Systematic, not accidental --- shared style class.
\end{knote}

\subsection*{The self-test}

\walkbeat{The student scrolls to the footer and presses Run self-tests. The panel fills in a single frame.}

\walkfig{s1-tests}{The suite, run live: 114 of 114 passing.}

\Stu{114 of 114, and the button relabels itself with the split: 62 computational, 52 rendered-text pins. Two species, and the count says which is which because you cannot judge the suite without knowing the mixture.}

\Kle{Define the second species.}

\Stu{A rendered-text pin reads the actual DOM after render and asserts that the prose says what it is supposed to say. The Q and A one, for instance, asserts there are exactly 14 question items, that all 14 are closed at load, that at least ten jump buttons exist, and that the words ``Mazur,'' ``Artin--Verdier,'' ``Chebotarev'' and ``It is not a functor'' are present in the rendered text.}

\Kle{Then be precise about what that species can and cannot catch. It can catch the prose drifting away from the code. It cannot catch the prose and the code being wrong together, and it cannot catch a theorem being misattributed --- if you cited Milnor 1962 for something Turaev proved, all 52 would stay green.}

\Stu{That is exactly right and it is the limit of the whole suite. The pins catch drift. They do not catch error. What the 62 computational ones catch is different: those assert against expected values that do not come from the same run.}

\Kle{Give me one I can check sitting here.}

\Stu{Line six. \texttt{tonelliShanks(13,937) === [386,551]}.}

\walkbeat{Klein does the arithmetic on the back of the handout.}

\Kle{$386^2 = 148996$. And $937 \times 159 = 148983$. Difference 13. So $386^2 \equiv 13 \pmod{937}$, and the other root has to be $937 - 386 = 551$. Your square root is a square root.}

\Stu{And that is the point of that species: you did not have to trust the implementation, you had to trust multiplication.}

\Kle{Which is the correct place to put the trust. Now --- what happens if I break one?}

\Stu{The failing line turns red and bold, and the button text changes to the passing count out of 114 followed by ``see failures.'' It does not tell you which line, so you scroll.}

\Kle{Through 114 monospace lines in no particular order, with no grouping by tab. On a green run that does not matter. On a red run it is the only moment the suite is useful and it is the moment it helps least.}

\Stu{Fair. There is no grouping and no jump-to-failure.}

\walkbeat{Klein runs his finger down the list and stops partway.}

\Kle{``Klein F1 pure.'' ``Klein A2 pure.'' ``Klein A4 pure.'' ``DOM Klein A1/A2.'' ``Klein C8 pure.''}

\Stu{Those are from your earlier review rounds. Each objection you made that was specific enough to be falsifiable got turned into a test with your initials on it, so it cannot silently come back.}

\Kle{A4 is the Chebotarev tally.}

\Stu{427 unramified odd primes below 3000. Identity 50, corner 47, and the three character classes 109, 111, 110 --- the observed frequencies hugging one eighth, one eighth, one quarter, one quarter, one quarter. That is the finite shadow of the density statement the Q and A tab makes in prose, pinned as a number.}

\Kle{And the certificate one, four from the bottom.}

\Stu{$[13,61,269] = -1$ via the certificate $(23,6,1)$, and then the $k=2$ scaling $(46,12,2)$ --- which satisfies the same identity --- would report $+1$, because scaling by $k$ multiplies the value by $(k/q)$ and $(2/269) = -1$. The engine rejects it. That test exists because the symbol is only well defined on normalized certificates, and the failure mode is silent otherwise.}

\Kle{So the suite's most interesting entries are all of the same kind: not ``the code gets the right answer,'' but ``the code refuses the wrong question.'' The rejection tests are worth more than the computation tests, because a wrong number is visible and a silently accepted invalid input is not.}

\Stu{The rejection taxonomy in \texttt{redeiSymbol} was the thing that took longest to get right. Primality, definedness mod 4, the identity, primitivity, $y$ even, $x - y \equiv 1 \pmod 4$, and then the case where $q$ divides a conjugate --- each with its own message, in that order, so a bad input hears the real reason first.}

\Kle{Then here is my summary of the suite, and it is a mixed one. 114 green is not a proof of anything and you have not claimed it is --- the split in the button is you telling me so before I ask, which I note. What it is proof of is that this app is a piece of engineering with a regression net under it, which is not the normal condition of a teaching demo, and which is why I have been able to spend two hours arguing about mathematics instead of about whether the numbers are stable.}

\begin{knote}
Ran the suite live: 114/114, split shown on the button as 62 computational, 52 rendered-text pins. Naming the split unprompted is the right instinct --- I did not have to ask what fraction were tautologies.

Spot-checked \texttt{tonelliShanks(13,937) = [386,551]} by hand: $386^2 = 148996 = 937 \cdot 159 + 13$, and $551 = 937 - 386$. Passes.

Limits, which he stated before I did: the 52 pins catch prose/code drift, not error --- a misattributed theorem stays green. The 62 computational assert against constants chosen off-run, so they are checkable, and I checked one.

Best entries are the rejection tests, not the computation tests: the scaled certificate $(46,12,2)$ satisfying the same identity and reporting $+1$ via $(2/269) = -1$, refused. A wrong number is visible; a silently accepted invalid input is not.

Weakness: 114 flat monospace lines, no grouping, no jump-to-failure. Useless exactly when it matters.

My initials are in five test names. That is either flattery or good engineering practice and I suspect it is both.
\end{knote}